% Options for packages loaded elsewhere
\PassOptionsToPackage{unicode}{hyperref}
\PassOptionsToPackage{hyphens}{url}
\PassOptionsToPackage{dvipsnames,svgnames,x11names}{xcolor}
%
\documentclass[
  10pt,
  a4paper,
]{article}
\usepackage{amsmath,amssymb}
\usepackage{setspace}
\usepackage{iftex}
\ifPDFTeX
  \usepackage[T1]{fontenc}
  \usepackage[utf8]{inputenc}
  \usepackage{textcomp} % provide euro and other symbols
\else % if luatex or xetex
  \usepackage{unicode-math} % this also loads fontspec
  \defaultfontfeatures{Scale=MatchLowercase}
  \defaultfontfeatures[\rmfamily]{Ligatures=TeX,Scale=1}
\fi
\usepackage{lmodern}
\ifPDFTeX\else
  % xetex/luatex font selection
  \setmainfont[]{Liberation Serif}
  \setsansfont[]{Liberation Sans}
  \setmonofont[]{DejaVu Sans Mono}
\fi
% Use upquote if available, for straight quotes in verbatim environments
\IfFileExists{upquote.sty}{\usepackage{upquote}}{}
\IfFileExists{microtype.sty}{% use microtype if available
  \usepackage[]{microtype}
  \UseMicrotypeSet[protrusion]{basicmath} % disable protrusion for tt fonts
}{}
\makeatletter
\@ifundefined{KOMAClassName}{% if non-KOMA class
  \IfFileExists{parskip.sty}{%
    \usepackage{parskip}
  }{% else
    \setlength{\parindent}{0pt}
    \setlength{\parskip}{6pt plus 2pt minus 1pt}}
}{% if KOMA class
  \KOMAoptions{parskip=half}}
\makeatother
\usepackage{xcolor}
\usepackage[margin=22mm]{geometry}
\setlength{\emergencystretch}{3em} % prevent overfull lines
\providecommand{\tightlist}{%
  \setlength{\itemsep}{0pt}\setlength{\parskip}{0pt}}
\setcounter{secnumdepth}{-\maxdimen} % remove section numbering
\ifLuaTeX
\usepackage[bidi=basic]{babel}
\else
\usepackage[bidi=default]{babel}
\fi
\babelprovide[main,import]{korean}
\ifPDFTeX
\else
\babelfont{rm}[]{Liberation Serif}
\fi
% get rid of language-specific shorthands (see #6817):
\let\LanguageShortHands\languageshorthands
\def\languageshorthands#1{}
\usepackage{xeCJK}
\setCJKmainfont{Noto Serif CJK KR}
\setCJKsansfont{Noto Sans CJK KR}
\setCJKmonofont{Noto Sans Mono CJK KR}
\xeCJKsetup{CJKspace=true}
\usepackage{xurl}
\usepackage{fvextra}
\usepackage{fancyhdr}
\AtBeginDocument{\DefineVerbatimEnvironment{verbatim}{Verbatim}{breaklines=true,breakanywhere=true,fontsize=\small}}
\setlength{\emergencystretch}{4em}
\setlength{\headheight}{15pt}
\pagestyle{fancy}\fancyhf{}
\fancyhead[L]{\small 중성자별 구조와 비방사 진동}
\fancyfoot[C]{\thepage}
\renewcommand{\headrulewidth}{0.3pt}
\setcounter{tocdepth}{2}
\ifLuaTeX
  \usepackage{selnolig}  % disable illegal ligatures
\fi
\usepackage{bookmark}
\IfFileExists{xurl.sty}{\usepackage{xurl}}{} % add URL line breaks if available
\urlstyle{same}
\hypersetup{
  pdftitle={A09 · Cowling single-shot와 표면 잔차},
  pdflang={ko-KR},
  colorlinks=true,
  linkcolor={Maroon},
  filecolor={Maroon},
  citecolor={Blue},
  urlcolor={Blue},
  pdfcreator={LaTeX via pandoc}}

\title{A09 · Cowling single-shot와 표면 잔차}
\usepackage{etoolbox}
\makeatletter
\providecommand{\subtitle}[1]{% add subtitle to \maketitle
  \apptocmd{\@title}{\par {\large #1 \par}}{}{}
}
\makeatother
\subtitle{과제 큐시트}
\author{}
\date{}

\begin{document}
\maketitle

\setstretch{1.12}
\textbf{선행:} A06, A08　 \textbf{독립 작업:} 7--12시간 (계획치)

\subsection{한 번 끝까지 적분한 해는 이미 진동
모드인가?}\label{uxd55c-uxbc88-uxb05duxae4cuxc9c0-uxc801uxbd84uxd55c-uxd574uxb294-uxc774uxbbf8-uxc9c4uxb3d9-uxbaa8uxb4dcuxc778uxac00}

유체 요소를 옮겼을 때 원래 위치에서 본 압력 변화와 요소를 따라간 압력
변화는 같을까? 표면에서 0으로 두어야 할 쪽을 먼저 말하자.

먼저 한두 문장으로 답한 뒤 아래 설명을 읽자. 첫 예측이 틀렸거나 아직
모르겠다고 적는 것은 감점하지 않는다.

\subsection{필요한 관계와 읽을
위치}\label{uxd544uxc694uxd55c-uxad00uxacc4uxc640-uxc77duxc744-uxc704uxce58}

\[
\begin{aligned}
Q&=(\omega/c)^2r^2e^{\lambda-2\nu}V+\nu'W,\\
W'&=Q/\chi-\ell(\ell+1)e^\lambda V,\\
V'&=2\nu'V-e^\lambda W/r^2,\\
\Delta P&=-(\epsilon+P)e^{-\lambda}Q/r^2.
\end{aligned}
\]

\begin{itemize}
\tightlist
\item
  \path|notes/chapters/ch05.md|
\item
  \path|notes/equations/cowling.md|
\item
  \path|learner/src/modes/ns_cowling.f90|
\item
  \path|learner/src/modes/ns_center_state.f90|
\item
  \path|learner/src/modes/ns_surface_residual.f90|
\end{itemize}

\subsection{계산으로
옮기기}\label{uxacc4uxc0b0uxc73cuxb85c-uxc62euxae30uxae30}

TOV 경험자는 제5장의 유도와 한 점 검산부터 시작한다. 중심·우변·표면
함수를 직접 작성할 수 있지만, 적응 적분과 양방향 연결은 제공 방법을
사용한다.

\begin{enumerate}
\def\labelenumi{\arabic{enumi}.}
\tightlist
\item
  제5장의 풀이를 따라 W,V의 차원과 중심 정칙해를 확인한다.
\item
  한 배경점 우변을 손계산하고 calculation의 출력 또는 자신의 물리 함수와
  대조한다.
\item
  한 trial에서 잔차의 진폭·방향 변환을 확인한다. 제공 양방향 matching은
  다시 구현하지 않는다.
\end{enumerate}

\newpage

\subsection{실행해서 무엇을 볼
것인가}\label{uxc2e4uxd589uxd574uxc11c-uxbb34uxc5c7uxc744-uxbcfc-uxac83uxc778uxac00}

패키지 루트에서 실행한다. \texttt{-\/-plan}은 경로만 확인하며 계산
성공을 뜻하지 않는다.

\begin{verbatim}
python tools/run_assignment.py A09 --plan
python tools/run_assignment.py A09 --run \
  --workspace learner --output ../A09_result.json
\end{verbatim}

\subsection{결과를
설명하기}\label{uxacb0uxacfcuxb97c-uxc124uxba85uxd558uxae30}

표면조건의 주체는 움직이는 유체 요소다. 한 점의 우변, 중심의 정칙성, 한
trial의 잔차를 확인해도 고유값 탐색과 모드 식별은 아직 남는다.

\textbf{남길 증거.} 차원표와 point RHS; 한 trial의 W,V profile; 잔차
진폭 불변성과 이번 값이 아직 고유모드가 아닌 이유; 같은 표면조건을
만족하는 양방향 시작 벡터와 matching 조건의 설명

\textbf{판단 기준.} 2점 예: 경계 잔차가 작아도 mode identification은
별도라는 점을 실행 상태와 연결한다. 최초 예측의 정오보다, 결과를 본 뒤
이유를 고치고 다른 입력에도 적용하는지를 본다.

\subsection{막혔을 때}\label{uxb9c9uxd614uxc744-uxb54c}

\textbf{힌트 1.} Q 각 항의 차원을 먼저 확인하고 부호는 별도 점검한다.

\textbf{힌트 2.} 양의 진폭 배율에서는 잔차가 유지되고 전체 부호
반전에서는 잔차 부호만 바뀌어야 한다.

\subsection{다음 만남에 다시
묻기}\label{uxb2e4uxc74c-uxb9ccuxb0a8uxc5d0-uxb2e4uxc2dc-uxbb3buxae30}

다음 회차에는 lapse에 상수를 더한 경우 주파수를 어떻게 바꿔야 같은 항을
얻는지 설명한다. 차원 검사만으로 찾지 못하는 오류도 하나 든다. 먼저 짧게
답한 뒤 본문과 자신의 기록을 열어 확인한다.

\textbf{범위.} 단일 trial 성공은 mode identity나 spectral completeness가
아니다. 평형 barotropic χ는 frozen composition을 포함하지 않는다.

\end{document}
